\subsection{Attribute-Based Encryption}
Attribute-Based Encryption, as the name suggest, is a type of encryption that enforce an access control via attributes and authorization policies. The secret key of a user and the ciphertext are both dependent by attributes.

Proposed in 1984, it was improved and implemented in the first years of 2000s; it presents two main types of encription schemes: the \textbf{key-policy attribute-based encryption (KP-ABE)} \cite{goyal2006attribute} and the \textbf{ciphertext-policy attribute-based encryption (CP-ABE)} \cite{bethencourt2007ciphertext}.

\subsubsection{From Identity-Based to Attribute-Based Encryption}
Identity-Based Encryption is a type of public key cryptography based on PKI (Public-Key Infrastructure), with the main difference being the fact that IBE allows a 1 to 1 relationship between sender and reciever.
As \citeauthor{shamir1984identity} said in his paper \cite{shamir1984identity}:
\begin{quote}
    The scheme is based on a public key cryptosystem with an extra twist: Instead of generating a random pair of public/secret keys and
publishing one of these keys, the user chooses his name and network address
as his public key. Any combination of name, social security number, street address, office number or telephone number can be used (depending
on the context) provided that it uniquely identifies the user
in a way he cannot later deny, and that it is readily available to the other party.
\end{quote}
The signature scheme abide the following verification condition \cite{shamir1984identity}:
\begin{equation}
    s^e = i \cdot t^{f(t,m)} \mod{n}
\end{equation}
Where 
\begin{itemize}
    \item $m$ is the message,
    \item $s,t$ is the signature,
    \item $i$ is the user's identity,
    \item $n$ is the product of large primes,
    \item $e$ is a large prime reltively prime to $\phi(n)$,
    \item $f$ is a one way function.
\end{itemize}

Its application are quite immediate to understand, like the ability to ensure that only a given person could read a message, without relying on the PKI is quite remarcable. There is, anyway, a trust assumtpion towards the Private Key Generator (PKG), which if compromised would nullify the secrecy of every encrypted data that use its issued keys.

The same assumption holds for the Attribute-Based Encryption, but now we have a broader use case: since both ciphertext and keys aren't associated with a single identity it is possible for the encryptor to avoid knowing the specific person who has to read, but, instead, the proprierties that a user has to have.

\subsubsection{Ciphertext-Policy versus Key-Policy ABE}
\subsubsection{Multi-authority ABE}
\subsubsection{Functional Encryption as a Generalization}

\subsection{Identity and Access Management}
\subsubsection{OAuth 2.0 and PKCE}
\subsubsection{Keycloak}

\subsection{The Model Context Protocol}

\subsection{Hook-Based Interception in AI Coding Agents}
